
This section introduces the theoretical framework within which we will
analyze the \bestofnalg and \rlhfalg algorithms. We first introduce
the \mlsharp objective as a stylized self-reward function, then
introduce our statistical framework for sharpening.
\iclr{We write $f=\bigoht(g)$ to denote 
    $f = \bigoh(g\cdot{}\max\crl*{1,\mathrm{polylog}(g)})$ and
    $a\approxleq{}b$ as shorthand for $a=\bigoh(b)$.\loose}
\subsection{Maximum-Likelihood Sharpening}\label{sec:autoreg_sharpening}

Our theoretical results focus on the \mlsharp objective given
by
\begin{align}
  \label{eq:ml_sharpening2}
    \rself(y \mid x) := \log \piref(y \mid x),
\end{align}
which we aim to maximize using conditional samples
$y\sim\piref(\cdot\mid{}x)$ from the base model. This is a simple and stylized self-reward function, but we will show
that it enjoys a rich theory. In
particular, we can restate the problem of sharpening with this self-reward through the lens of \emph{amortization}.
\loose
\begin{mainbox}
  \emph{Can we efficiently \textbf{amortize maximum likelihood
      inference (optimization)}
  for a conditional distribution $\piref(y\mid{}x)$ given access to a
  \textbf{sampling oracle} that can sample $y\sim\piref(\cdot \mid{}x)$?\loose}
\end{mainbox}
\dfedit{The tacit assumption in this framing is that the
  maximum-likelihood response constitutes a useful form of hidden
  knowledge.} \Mlsharp connects
the study of self-improvement to a large body of research in
theoretical computer science demonstrating computational reductions between
optimization (inference) and sampling (generation)
\citep{kirkpatrick1983optimization,lovasz2006fast,singh2014entropy,ma2019sampling,talwar2019computational}. %
Our sharpening framework offers a new learning-theoretic perspective by focusing on the problem of amortizing this
type of reduction\akdelete{ \abedit{and is inspired by the classical query complexity framework\citep{nemirovski1983problem,traub1988information,raginsky2011information,agarwal2012information}}}. \msdelete{}

We evaluate the quality of an
approximately sharpened model as follows. Let
\[\Ystar(x)\ldef{}\argmax_{y\in\cY}\log \piref(y\mid{}x);\] we interpret
$\Ystar(x)\subset\cY$ as a set  to accommodate non-unique
maximizers, and will write $\ystar(x)$ to indicate a
unique maximizer when it exists (i.e., when
$\Ystar(x)=\crl{\ystar(x)}$). 


                                                     \begin{definition}[Sharpened model]
                                                       \label{def:sharpening}
    We say that a model $\pihat$ is $(\eps,\delta)$-sharpened
    relative to $\piref$ if
    \begin{align}
      \bbP_{x\sim\cdist}\brk*{\pihat\prn*{\Ystar(x)\mid{}x}\geq{}1-\delta} \geq{} 1-\eps.
    \end{align}
  \end{definition}
That is, an $(\eps,\delta)$-sharpened model places at least
$1-\delta$ mass on arg-max responses on all but an $\eps$-fraction
of prompts under $\mu$. For small $\delta$ and $\eps$, we are
guaranteed that $\pihat$ is a high-quality generator: sampling from the model will produce an arg-max
response with high probability for most prompts.







  \paragraph{\Mlsharp for autoregressive models}
  Though our most general results are
  agnostic to the structure of $\cX$, $\cY$, and $\piref$, our primary
  motivation is
the autoregressive setting in which
$\cY=\cV^{H}$ for a \emph{vocabulary space} $\cV$ and sequence
  length $H$, and where $\piref$ has the autoregressive structure
$\piref(y_{1:H}\mid{}x)=\prod_{h=1}^{H}\pirefh(y_h\mid{}y_{1:h-1},x)$
for $y=y_{1:H}\in\cY$.
We observe
that when the response
$y=(y_1,\ldots,y_H)\in\cY=\cV^{H}$ is a sequence of tokens, the
\mlsharp objective \eqref{eq:ml_sharpening} sharpens
toward the \emph{sequence-level} arg-max response:
\begin{align}
  \label{eq:inference_time}
\argmax_{y_{1:H}}\log\piref(y_{1:H}\mid{}x).
\end{align}
\dfedit{Although somewhat stylized, \cref{eq:inference_time}
 is a non-trivial (in general, computationally intractable; see \Cref{sec:hardness})
 solution concept. We view the sequence-level arg-max as a form of hidden
 knowledge that cannot necessarily be uncovered through naive
 sampling or greedy decoding.
  \loose}


  


  \loose

  \paragraph{Role of $\delta$ for autoregressive models}


  As can be verified through simple examples, beam-search and greedy
  tokenwise decoding do not return an exact (or even approximate)
  solution to \eqref{eq:inference_time} in general. There is one
  notable exception: If the model has already been sharpened to
  $\delta<1/2$ and the arg-max sequence is unique, then greedy
  decoding will succeed.
  \loose
  \begin{proposition}[Greedy decoding succeeds for sharpened policies]
    \label{prop:greedy}
    Let $\pi=\pi_{1:H}$ be an autoregressive model defined over
    response space $\cY=\cV^{H}$. For a given prompt $x\in\cX$, if
    $\Ystar(x)=\crl*{\ystar(x)}$ is a singleton and
    $\pi(\ystar(x)\mid{}x)>1/2$, then the greedy decoding strategy
    that selects
    \iclr{$\yhat_h=\argmax_{y_h\in\cV}\pi_h(y_h\mid{}\yhat_1,\ldots,\yhat_{h-1},x)$}
    \arxiv{\begin{align}
      \yhat_h=\argmax_{y_h\in\cV}\pi_h(y_h\mid{}\yhat_1,\ldots,\yhat_{h-1},x)
    \end{align}}
    guarantees that $\yhat=\ystar(x)$. This result is tight, in the
    sense that there exist $\pi$ with $\pi(\ystar(x)\mid{}x)\leq{}1/2$
    for which greedy decoding fails to recover $\ystar(x)$. \loose
  \end{proposition}
  This means that if we start from an un-sharpened model, it can suffice
  to focus on sharpening to $\delta<1/2$.








  






 

















